\documentclass[t,12pt]{beamer}

\usepackage{amssymb}
\usepackage{amsfonts}
\usepackage{amsmath}
\usepackage{dsfont}
\usepackage[utf8x]{inputenc}
\usepackage[labelfont=bf,justification=justified]{caption}
\usepackage{subcaption}
\usepackage{fixltx2e}
\usepackage{setspace}
\usepackage{arydshln}
\usepackage{xcolor}
\usepackage{tikz}
\usepackage{chngpage}
\usepackage{comment}
\usetikzlibrary{shapes,arrows}
\tikzstyle{block} = [rectangle, draw, fill=blue!20, 
    text width=5em, text centered, rounded corners, minimum height=4em]
\tikzstyle{line} = [draw, -latex']

\usetheme{Emory}
\makeatletter
\setbeamertemplate{footline}
{
  \leavevmode%
  \hbox{%
  \begin{beamercolorbox}[wd=.333333\paperwidth,ht=2.25ex,dp=1ex,center]{author in head/foot}%
    \usebeamerfont{author in head/foot}\insertshortauthor~~\beamer@ifempty{\insertshortinstitute}{}{\insertshortinstitute}
  \end{beamercolorbox}%
  \begin{beamercolorbox}[wd=.333333\paperwidth,ht=2.25ex,dp=1ex,center]{title in head/foot}%
    \usebeamerfont{title in head/foot}\insertshorttitle
  \end{beamercolorbox}%
  \begin{beamercolorbox}[wd=.333333\paperwidth,ht=2.25ex,dp=1ex,right]{date in head/foot}%
    \usebeamerfont{date in head/foot}\insertshortdate{}\hspace*{2em}
    \insertframenumber{} / \inserttotalframenumber\hspace*{2ex} 
  \end{beamercolorbox}}%
  \vskip0pt%
}
\makeatother

\setbeamerfont{frametitle}{size=\large}
\setbeamertemplate{navigation symbols}{}

\titlegraphic{\includegraphics[width=0.3\textwidth]{/home/william/Emory.png}}

\begin{document}

\title[AR(1) and unit-root processes]{AR(1) and unit-root processes}
\author{\vspace{0.35in} \\ {\fontfamily{ppl}\selectfont William Mann} \vspace{-0.35in}}

\date{}
\maketitle

\begin{frame}[c]{AR(1) process and random-walk process}
AR(1) means ``autoregressive process of order 1''. It looks like this:
$$
y_t \quad = \quad \mu \quad + \quad \phi y_{t-1} \quad + \quad u_t
$$
where each $u_t$ is an independent draw from $\text{N}(0,\sigma^2)$.
\\ \smallskip
An AR(1) process eventually stabilizes around an average, $\frac{\mu}{1-\phi}$. %
\\ \medskip
If $\phi=1$, we get a \textbf{random walk}, an example of a \textbf{unit-root} process:
$$
y_t \quad = \quad \mu \quad + \quad y_{t-1} \quad + \quad u_t
$$
This process does \emph{not} have any long-run ``average.''
\begin{itemize}
\item This is a very common case with financial data.
\item It also creates some challenges for statistical modeling!
\\ To understand why, we need the idea of a ``stationary'' process...
\end{itemize}
\end{frame}

\begin{frame}[c]{Weakly stationary process (or covariance stationary)}
\label{weakly_stationary}
A process $y$ is \textbf{weakly stationary} (or \textbf{covariance stationary}) if
\begin{itemize}
\medskip
\item $\mathbb{E}[y_t]$ does not depend on $t$,
\item $Var(y_t)$ does not depend on $t$,
\item $\text{Cov}(y_t,y_s)$ depends only on the distance $t-s$.
\end{itemize}
\bigskip
In words, this means that $y$ follows a stable distribution at all times.
\\ \vfill
{\footnotesize
From now on, I will use ``stationary'' as a shorthand for ``weakly stationary,''
\\
and will use ``nonstationary'' to mean ``not weakly stationary,''
\\ 
Other fields outside economics sometimes use these terms differently.
}
\end{frame}

\begin{frame}{Stationary and nonstationary examples}
An AR(1) with $\phi<1$ is stationary:  
\begin{align*}
\mathbb{E}[y_t] \quad &= \quad \frac{\mu}{1-\phi}
\\
Var(y_t) \quad &= \quad \frac{\sigma^2}{1-\phi^2}
\\
Cov(y_t,y_s) \quad &= \quad \phi^{t-s} \times \frac{\sigma^2}{1-\phi^2}
\end{align*}
A random walk is \emph{not} stationary:
\begin{align*}
\mathbb{E}[y_t] \quad &= \quad t \times \mu \quad + \quad y_0
\\
Var(y_t) \quad &= \quad t \times \sigma^2
\\
\text{Cov}(y_t,y_s) \quad &= \quad \min(t,s) \times \sigma^2
\end{align*}
\end{frame}

\begin{frame}[c]{The difference of a random walk is stationary}
\label{difference_stationary}
The \emph{difference} of a random walk \emph{is} stationary.
\\ \medskip
To see this, use $\Delta y_t$ to represent $y_t - y_{t-1}$. Then we can write
$$
\Delta y_t \quad = \quad \mu + u_t
$$
and from this we can see that
\begin{align*}
\mathbb{E}[\Delta y_t] \quad &= \quad \mu
\\
Var(\Delta y_t) \quad &= \quad \sigma^2
\\
\text{Cov}(\Delta y_t, \Delta y_s) \quad &= \quad 0
\end{align*}
This is an important fact that we will use a lot.
\end{frame}

\begin{frame}[c]{Why is stationary data important?}
\begin{itemize}
\item Statistics is all about \emph{estimating} averages, variances, etc.
\\ This implicitly assumes that such numbers actually exist!
\medskip
\item For a nonstationary process, they do not.
\\ So the results of such analysis would not be meaningful.
\\ We have to find some other approach instead.
\medskip
\item 
What happens if we (mistakenly) ignore this issue,
\\ and perform our usual analysis using nonstationary data?
\medskip
\item The computer will not stop us.
\\ It will still run the analysis and give us (meaningless) results.
\medskip
\item Instead, we have to check for nonstationary data ourselves, and figure out what to do about it.
This is our next topic.
\end{itemize}
\end{frame}

\begin{frame}[c]
\vfill
\centerline{Testing for unit roots}
\vfill
\end{frame}

\begin{frame}[c]{Types of nonstationary process}
\begin{itemize}
\item Statistical tests require a precise null and alternative hypothesis.
\medskip
\item In theory, there are many ways a process can be nonstationary.
\\ Each type of nonstationary behavior requires its own test.
\medskip
\item
In practice, the most plausible and common category of nonstationary processes is called \emph{unit root} processes.
\medskip
\begin{itemize}
\item
A unit-root process is one that is not stationary on its own,
\\ but can be turned into a stationary process through differencing.
\item As we saw earlier, a random walk process is one example.
\end{itemize}
\medskip
\item So we will focus our attention on tests for unit-root behavior.
\end{itemize}
\end{frame}

\begin{frame}[c]{Estimating AR(1) and AR(p) processes}
Unit root tests build on the following information:
\\ \bigskip
Recall the definition of an AR(1) process,
$$
y_t \quad = \quad \mu \quad + \quad \phi y_{t-1} \quad + \quad u_t
$$
Suppose you believe a data series follows an AR(1) process. 
\\ How would you estimate $\phi$? Just regress $y_t$ on $y_{t-1}$.
\\ \bigskip
An AR(p) process is like an AR(1), with $p$ lags of $y$:
$$
y_t \, = \, \mu \, + \, \phi_1 y_{t-1} + \phi_2 y_{t-2} + ... + \phi_p y_{t-p} \, + \, u_t
$$
You can estimate all the $\phi$ at once, by running a multiple regression of $y_t$ on $y_{t-1},\dots,y_{t-p}$.
See notebook on Canvas for examples.
\end{frame}

\begin{frame}[c]{Unit root tests: Dickey-Fuller}
The Dickey-Fuller unit root test assumes the data are AR(1), 
\\ and tests the null $\phi=1$ against the alternative $\phi < 1$.
\begin{itemize}
\item Write out $\Delta y_t$, which is stationary under either hypothesis:
$$
\Delta y_t \quad = \quad \mu \quad + \quad (\phi - 1) \times y_{t-1} \quad + \quad u_t
$$
\vspace{-0.2in}
\item Estimate $\phi-1$ by regressing $\Delta y_t$ on $y_{t-1}$.
\\ Then we want to test whether this coefficient is negative.
\\ If the $p$-value is low enough, we can reject that $\phi=1$.
\item The test statistic is the $t$-stat from the regression.
\\ But, the typical regression p-values will not be correct:
\\ Sampling distribution is nonstandard under the null $\phi=1$.
\item Refer to special tables of critical values to check significance.
\end{itemize}
\end{frame}

\begin{frame}[c]{Unit root tests: Augmented Dickey-Fuller (ADF)}
Same approach, but control for lags of $\Delta y_t$ in the regression,
\\ allowing the process to be any AR(p) instead of just AR(1).
\\ \medskip
For example, we can allow $y$ to be AR(2) by running
$$
\Delta y_t \quad = \quad \mu \quad + \quad (\phi - 1) \times y_{t-1} \quad + \quad \delta \times \Delta y_{t-1} \quad + \quad u_t
$$
Each extra lag allows more general dynamics in $y$.
\begin{itemize}
\item How far to go with this? There is no simple rule. 
\item One common procedure is to select based on information criteria such as AIC or BIC (see definitions next week).
\item Most software allows you to automate this process.
\end{itemize}
\medskip
In practice, ADF is the most common unit-root test.
\end{frame}

\begin{frame}[c]{Deterministic trend as alternative hypothesis}
\vspace{-0.15in}
In the tests that we described just now:
\begin{itemize}
\item the null hypothesis was that the process has a unit root, 
\item the alternative was that the process is stationary.
\end{itemize}
\medskip
Another potential behavior is a \emph{deterministic} trend.
\begin{itemize}
\item This is a trend that is a simple function of time.
\item Also called \emph{trend-stationary} behavior.
\end{itemize}
\medskip
This can be specified as the alternative in an ADF test by
\begin{itemize}
\item adding a time trend in the regression,
\item or, running the test on detrended data.
\end{itemize}
(In practice, just specify the appropriate option in software.)
\\ \medskip
However, deterministic trends are seen as being rare in practice.
\end{frame}

\begin{frame}[c]{Unit root as the alternative: KPSS}
The KPSS test reverses the ADF procedure:
\\ Stationary data is the null, unit root is the alternative.
\\ Hence we hope to \emph{not} reject the null under this test.
\\ \bigskip
The derivation is a bit too involved to describe here,
\\ but it's implemented in any statistical software.
\\ \bigskip
KPSS and ADF are often used together.
\\ We hope they give consistent answers (one reject, one fail to reject).
\\ But there is no guarantee of this. Both can reject, or fail to reject.
\\ \bigskip\bigskip
{\footnotesize (KPSS $=$ Kwiatkowski, Phillips, Schmidt, and Shin.)}
\end{frame}

\begin{frame}[c]{General comments about unit root tests}
\begin{itemize}
\item Statistical tests might give you the impression that they are definitive and conclusive, but they are not. No test is perfect. 
\medskip
\item Unit root tests in particular have low power. 
\\ This means they often struggle to give us any clear conclusion.
\medskip
\item By the same token, these tests do not give us a cookbook that we can follow to tell us exactly how to handle a time series.
\medskip
\item Instead, view them as a way to check intuition that you already have from background knowledge or economic theory. 
\end{itemize}
\end{frame}

\begin{frame}[c]{Deal with unit-root data by differencing}
The standard thing to do with unit-root data is to difference it.
\begin{itemize}
\smallskip
\item After one or more differencing steps, the data will be stationary.
\\ (In all our examples, one step will be enough.)
\item Statistical analysis on this differenced data will be valid.
\end{itemize}
\medskip
For example, instead of modeling the price of a stock or an index, 
\\ we model its \emph{returns}, which are the difference of log price.
\\ \medskip
The major sacrifice is that you give up on modeling $y_t$ directly, 
\\ and only attempt to explain the change $\Delta y_t$.
\\ \medskip
So, you don't want to difference data unless absolutely necessary.
\\ \medskip
But if the series has a unit root, you typically have no choice.
\end{frame}

\end{document}